% \vspace{-0.25em}
\subsection{Continual Learning and Self-Evolving Agents}
% \vspace{-0.25em}
Continual learning (CL) aims to enable models to learn sequentially while mitigating catastrophic forgetting~\citep{parisi2019continual,wang2024comprehensivesurveycontinuallearning}. While various replay-based and regularization methods have been proposed, most CL paradigms assume predefined task boundaries and focus on knowledge preservation rather than active acquisition in open-ended environments~\citep{Kirkpatrick_2017,ding2024boostinglargelanguagemodels,huai2025taskcorememorymanagementconsolidation,huai2025clmoeenhancingmultimodallarge}. The pursuit of self-evolving agents moves beyond these limitations by enabling systems to grow autonomously from experience. Frameworks such as Reflexion and Generative Agents explore self-improvement through self-play and reflective reasoning, often storing past trajectories as memory to guide future actions~\citep{shinn2023reflexionlanguageagentsverbal,wei2022chain,yao2023tree,besta2024graph,yao2023reactsynergizingreasoningacting}. However, these systems either store raw, unstructured data or rely on memory mechanisms that are not designed for the systematic, long-term distillation and refinement of abstract strategic knowledge. 
% Our EvolveR framework introduces a novel, closed-loop paradigm for self-evolution. 
Instead of relying on external data streams, our agent autonomously generates and refines its own experiences through an iterative cycle of online interaction and offline reflection.


\subsection{LLM Agents and Reinforcement Learning}
LLM agents have been widely explored through frameworks such as ReAct, which interleaves reasoning and actions, and Reflecion, which improves task performance via self-reflection~\citep{yao2023reactsynergizingreasoningacting,shinn2023reflexionlanguageagentsverbal}. While these approaches are primarily prompt-based and stateless, they prevent long-term accumulation of strategic knowledge. External memory frameworks like ExpeL address this limitation by reusing past trajectories, but they do not enable systematic self-improvement across tasks~\citep{zhao2024expelllmagentsexperiential}. While effective, these methods often rely on simple prompting and are inherently stateless, limiting their ability to internalize knowledge across tasks. Recent work has increasingly turned to reinforcement learning (RL) to train agents for long-horizon, multi-turn tasks. However, applying RL is challenging due to sparse rewards and the need for stable training signals. Search-R1~\citep{jin2025searchr1trainingllmsreason}, O2-Searcher~\citep{mei2025o2searchersearchingbasedagentmodel}, and AutoRefine~\citep{shi2025searchrefinethinkautonomous} all use RL to train LLMs to generate and interact with external search tools. While these works successfully optimize the LLM's interaction with external factual knowledge, they do not address the broader challenge of an agent's self-improvement through its own internal experience.



